Four bodies of work ground the design decisions made in this thesis, each informing a different
part of the system.

\section{Natural Language Reasoning over Virtual Knowledge Graphs}
\label{sec:rel-nl2sparql}

Balasooriya, Dalla Vecchia, and Quintarelli~\cite{balasooriya2026} describe the baseline system
this thesis extends: Ontop exposes relational Verona tourism data as a VKG; NL2SPARQL translation
is delegated to SPARQL-LLM, an external, triplestore-agnostic, metadata-driven Text2SPARQL system
using SHACL-annotated query examples, VoID schema descriptions, and embedding-based retrieval; and
a GPT-4.1-mini prompt-based intent classifier first decides whether a question is factual (routed
to NL2SPARQL) or analytical/preference-oriented (routed to a lightweight, frequency-based inference
module over Verona Card visit data). Their preliminary evaluation reports 100\% routing accuracy on
50 hand-labeled questions, but only 70\% SPARQL-correctness on the factual subset --- 12 of 40
failures attributed to entity-linking errors or invalid syntax. Their conclusion explicitly names
recommendations as future work; this thesis is that future work, generalized to a second domain and
grounded in real interaction data rather than only aggregate frequency.

Three design implications follow directly. First, the existing pipeline's stage-by-stage structure
is kept intact; this thesis adds a new module operating \emph{after} result formatting, rather than
replacing the NL2SPARQL core (Chapter~\ref{ch:arch}). Second, their intent-routing pattern --- an
LLM classifying before acting --- is a template this thesis does not directly reuse, but their
70\% SPARQL-correctness figure is a direct, empirically-grounded warning: entity-linking and query
validity are real failure modes, and any downstream step that generates \emph{new} natural-language
queries inherits an analogous risk that they might not be answerable at all
(Section~\ref{sec:rec-validation}).

\section{LLM-Based Query Recommendation}
\label{sec:rel-gqr}

Bacciu, Palumbo, Damianou, Tonellotto, and Silvestri~\cite{bacciu2024} reframe query recommendation
as a pure generation task: prompt an LLM with a handful of \emph{query $\to$ recommendations}
example pairs plus the user's current query, and let the model generate recommendations directly,
with no query-log-trained ranking model needed --- solving the cold-start problem that afflicts
classical, log-mining recommenders. Their RA-GQR variant retrieves semantically similar past
queries from a log via sentence embeddings and approximate nearest-neighbor search, and uses those,
rather than generic hand-curated examples, to build the few-shot prompt dynamically --- Retrieval-
Augmented Generation applied to the recommendation prompt itself. In their experiments, RA-GQR
measurably outperforms plain GQR (roughly $+13$ to $+22\%$ NDCG@10).

This is the direct technique behind Section~\ref{sec:rec-retrieval} and~\ref{sec:rec-generation}: a
prompt is built from (i) a small number of retrieved, embedding-similar past queries from the
user's own log, (ii) the current query, and (iii) an instruction to continue with recommendations,
following their prompt template structure. Their evaluation protocols and metrics --- scoring a
recommendation as if it substituted for the original query, or as if it were appended to it, plus
a blind user-preference study --- are the direct ancestor of this thesis's own ablation design in
Chapter~\ref{ch:eval}, adapted to a query-history-versus-behavioral-profile axis in place of their
original data-availability axis.

\section{LLM and Knowledge Graph Integration}
\label{sec:rel-kgllm}

Pan, Luo, Wang, Chen, Wang, and Wu~\cite{pan2024} survey a broad taxonomy of LLM+KG integration
patterns. Two parts are directly relevant here, not the full taxonomy: their KG-prompting/KGQA
sections describe patterns for injecting retrieved graph facts into an LLM prompt so it reasons
\emph{with} the graph rather than purely from parametric knowledge, informing how this thesis's
recommended queries should ideally be grounded in the loaded ontology's actual classes and
properties rather than an imagined schema; and their repeated, explicit documentation of a
hallucination/faithfulness risk --- that LLM-generated content referencing structured knowledge
must match the target graph \emph{exactly}, and that this remains a known failure mode even with
in-context learning. Combined with the baseline system's own measured 70\% SPARQL-correctness rate
(Section~\ref{sec:rel-nl2sparql}), this directly motivates \emph{not} shipping raw LLM-suggested
queries to a user without first checking that they are actually answerable against the currently
loaded schema --- the requirement Section~\ref{sec:rec-validation} attempts, and partially fails,
to satisfy.

\section{Classical, Non-LLM Query Recommendation}
\label{sec:rel-qrmoccga}

Barman, Sarkar, and Chowdhury~\cite{barman2024} model query recommendation as multi-objective
optimization over a query-click bipartite graph built from search logs, using two objective
functions built from four concrete similarity features --- compound click probability, Jaccard
index of clicked links, Jaccard index of query words, and longest-common-subsequence ratio of query
strings --- optimized via a cooperative co-evolutionary genetic algorithm (QRMOCCGA), evaluated
against SimRank, Heat Diffusion, and commercial search engines on real AOL query logs.

Reimplementing their full cooperative co-evolutionary genetic algorithm is out of scope for this
thesis, but two of its ideas are directly useful and were borrowed at negligible cost. First, their
behavioral-similarity features --- word-overlap and click-overlap similarity between two queries ---
are a simple, non-LLM way to quantify ``how well does this candidate query match the user's actual
watch behavior,'' used directly inside this thesis's own ranking step
(Section~\ref{sec:rec-ranking}) without training anything. Second, their explicit two-objective
framing --- semantic relevance versus lexical/visual similarity --- is the direct conceptual
ancestor of this thesis's own two ranking axes, query-history similarity versus behavioral
alignment, even though the actual ranking mechanism here is a simple weighted score with greedy
diversification rather than a genetic algorithm. QRMOCCGA's simpler feature set is also used
directly as a classical, non-LLM baseline in this thesis's evaluation (Section~\ref{sec:eval-results}).

\section{Positioning of This Thesis}
\label{sec:rel-position}

No prior work reviewed here combines all three ingredients this thesis's recommender uses jointly:
retrieval-augmented, LLM-based candidate generation (from GQR/RA-GQR); an explicit, per-user
divergence signal between stated and actual preference, rather than either signal alone (motivated
directly by the supervisor's framing and absent from all four reviewed papers); and a
schema-grounding validation step informed by, but going beyond, generic LLM+KG hallucination
concerns, calibrated empirically against this project's own ontologies rather than assumed to work.
The negative empirical result reported in Section~\ref{sec:rec-validation} --- that class-level
embedding similarity cannot reliably separate on-topic from off-topic candidates when source
ontology metadata is this sparse --- is, to the author's knowledge, not previously documented in
this specific combination of techniques, and is offered as a concrete, reportable finding for
future work in this space to build on.
